\documentclass[final]{beamer}

\usepackage[scale=1.24]{beamerposter} % Use the beamerposter package for laying out the poster
\usepackage{bbm}
\usepackage{graphicx}
\usepackage{textpos}
\usepackage{xcolor}
\graphicspath{{./figs/}}
\usetheme{confposter} % Use the confposter theme supplied with this template

\setbeamercolor{block title}{fg=illiniblue,bg=white} % Colors of the block titles
\setbeamercolor{block body}{fg=black,bg=white} % Colors of the body of blocks
\setbeamercolor{block alerted title}{fg=white,bg=illiniorange} % Colors of the highlighted block titles
\setbeamercolor{block alerted body}{fg=black,bg=heritageorange} % Colors of the body of highlighted blocks
% Many more colors are available for use in beamerthemeconfposter.sty

%-----------------------------------------------------------
% Define the column widths and overall poster size
% To set effective sepwid, onecolwid and twocolwid values, first choose how many columns you want and how much separation you want between columns
% In this template, the separation width chosen is 0.024 of the paper width and a 4-column layout
% onecolwid should therefore be (1-(# of columns+1)*sepwid)/# of columns e.g. (1-(4+1)*0.024)/4 = 0.22
% Set twocolwid to be (2*onecolwid)+sepwid = 0.464
% Set threecolwid to be (3*onecolwid)+2*sepwid = 0.708

\newlength{\sepwid}
\newlength{\onecolwid}
\newlength{\twocolwid}
\newlength{\threecolwid}
\newlength{\fourcolwid}
\setlength{\paperwidth}{48in} % A0 width: 46.8in
\setlength{\paperheight}{36in} % A0 height: 33.1in
\setlength{\sepwid}{0.024\paperwidth} % Separation width (white space) between columns
\setlength{\onecolwid}{0.22\paperwidth} % Width of one column
\setlength{\twocolwid}{0.464\paperwidth} % Width of two columns
\setlength{\threecolwid}{0.708\paperwidth}
\setlength{\fourcolwid}{0.9508\paperwidth} % Width of three columns
\setlength{\topmargin}{-0.75in} % Reduce the top margin size
%-----------------------------------------------------------

\usepackage{graphicx}  % Required for including images
\usepackage[T1]{fontenc}
\usepackage{charter}
\usepackage[expert]{mathdesign}
\usepackage[percent]{overpic}
%\usepackage{caption}
\usepackage{booktabs} % Top and bottom rules for tables\usepackage[T1]{fontenc}
\usepackage{charter}

\DeclareFontShape{OMX}{cmex}{m}{n}{
  <-7.5> cmex7
  <7.5-8.5> cmex8
  <8.5-9.5> cmex9
  <9.5-> cmex10
}{}
\SetSymbolFont{largesymbols}{normal}{OMX}{cmex}{m}{n}
\SetSymbolFont{largesymbols}{bold}  {OMX}{cmex}{m}{n}


%----------------------------------------------------------------------------------------
%	TITLE SECTION 
%----------------------------------------------------------------------------------------

\title{\Large Crop Rotations, Transition Costs, and Farm Delinquencies} % Poster title

\author{Victor Funes-Leal, Lawson Connor and Eunchun Park} % Author(s)

\institute{University of Arkansas\\ Department of Agricultural Economics and Agribusiness} % Institution(s)

%----------------------------------------------------------------------------------------

\begin{document}
\addtobeamertemplate{headline}{} 
{
%\begin{tikzpicture}[remember picture,overlay] 
%\node [shift={(-10 cm,-7cm)}] at (current page.north east) %{
\includegraphics[scale=2]{BlockI-Logo}}; 
%\end{tikzpicture} 

\begin{tikzpicture}[remember picture,overlay] 
	\node [shift={(14 cm,-7cm)}] at (current page.north west) {
\includegraphics[scale=0.75]{figs/UA_Logo.eps}}; 
\end{tikzpicture} 
}

% {
% \begin{tikzpicture}[remember picture,overlay] 
% \node [shift={(-1 cm,-7cm)}] at (current page.north west) {\includegraphics[scale=.1]{bucky}}; 
% \end{tikzpicture} 
% }
\addtobeamertemplate{block end}{}{\vspace*{2ex}} % White space under blocks
\addtobeamertemplate{block alerted end}{}{\vspace*{2ex}} % White space under highlighted (alert) blocks

\setlength{\belowcaptionskip}{2ex} % White space under figures
\setlength\belowdisplayshortskip{2ex} % White space under equations

\begin{frame}[t] % The whole poster is enclosed in one beamer frame
\begin{columns}[t] % The whole poster consists of three major columns, the second of which is split into two columns twice - the [t] option aligns each column's content to the top

\begin{column}{\sepwid}\end{column} % Empty spacer column

\begin{column}{\onecolwid} % The first column

%----------------------------------------------------------------------------------------
%	OBJECTIVES
%----------------------------------------------------------------------------------------

\begin{block}{Motivation}
\begin{itemize}
    \item Operating credit is the financial lifeline of row-crop farms, yet credit-risk models treat revenue as one aggregate stream and ignore how rotation history shapes next year's yields.
    \item This study quantifies how three Illinois rotation plans differ in cash-flow delinquency risk (the borrower's problem) and peak credit demand (the lender's problem) over a six-year horizon.
\end{itemize}
\end{block}

\begin{block}{Research questions}
\begin{enumerate}
    \item Does a more complex rotation improve financial outcomes, or do agronomic gains hit diminishing returns versus simple alternation?
    \item What is the 6-year probability of at least one operating-loan delinquency, conditioning yields on the full rotation history?
    \item How does the negative price–yield correlation change those probabilities?
\end{enumerate}
\end{block}

\begin{block}{Data \& methods}
\begin{itemize}
\item \textbf{Yields}: QDANN yield-mapping model, every IL field 2008–2022, productivity defined by quantiles of this yield distribution.
\item \textbf{Rotations}: 6-year histories from the Cropland Data Layer (CDL).
\item \textbf{Weather/soil}: GRIDMET, TerraClimate (VPD, PDSI), gSSURGO.
\item \textbf{Costs}: 2024 Farmdoc Crop Budgets (vary by region and productivity in Central IL).
\item \textbf{Prices \& costs}: grain-elevator prices; copula embeds the negative price–yield correlation.
\item \textbf{Credit model}: spring origination $\rightarrow$ 9-month accrual $\rightarrow$ harvest repayment $\rightarrow$ balance.
rollover
\end{itemize}
\end{block}
\end{column} % End of the first column



\begin{column}{\sepwid}\end{column} % Empty spacer column

\begin{column}{\twocolwid} % Begin a column which is two columns wide (column 2)
	
\begin{columns}[t,totalwidth=\twocolwid] % Split up the two columns wide column

\begin{column}{\onecolwid}\vspace{-.6in} % The first column within column 2 (column 2.1)
\begin{block}{Rotation cuts 6-year delinquency risk by a third or more}
	%\vspace{-0.9in}
	\begin{figure}
		\begin{center}
			\includegraphics[scale=1.2]{figs/fig_p_bar_all.png}\\
				\caption{\hspace{.05cm}Probability of at least one operating-loan delinquency over the six-year horizon, by region (columns) and productivity zone (rows). Continuous corn (red) is highest everywhere; strict alternation (blue) is lowest.}
			\label{fig:fig_p_bar_all}
		\end{center}
	\end{figure} 
\end{block}
\end{column} % End of column 2.1

\begin{column}{\onecolwid}\vspace{-.6in} % The second column within column 2 (column 2.2)
\begin{block}{Survival curves trace where rotation pulls farms apart}
		%\vspace{-0.9in}
		\begin{figure}
			\begin{center}
				\includegraphics[scale=1.2]{figs/fig_p_surv_all.png}\\
				\caption{\hspace{.05cm}Pr(any missed payment) over the six-year horizon. The step structure mirrors the sawtooth: drops concentrate in corn years; soy years are flat. By year 6, alternation (blue) preserves $\approx$ 70–85\% of farms in solvency vs $approx$ 45–70\% under continuous corn (red).}
				\label{fig:fig_p_cf_all}
			\end{center}
		\end{figure} 
\end{block}
\end{column}	
\end{columns} % End of the split of column 2 - any content after this will now take up 2 columns width
	

\begin{columns}[t,totalwidth=\twocolwid] % Split up the two columns wide column

\begin{column}{\onecolwid}\vspace{-.6in} % The first column within column 2 (column 2.1)
\begin{block}{Peak loan demand stays under common credit limits}
	%\vspace{-0.9in}
	\begin{figure}
		\begin{center}
			\includegraphics[scale=1.2]{figs/fig_peak_loan_all.png}\\
			\caption{\hspace{.05cm}Distribution of peak gross spring loan (\$/ac) with \$700 / \$900 / \$1200 credit limits overlaid. Medians cluster near \$500–550; rotation shifts the whole distribution down, widening the margin to any binding ceiling.}
			\label{fig:fig_p_year_all}
		\end{center}
	\end{figure} 
\end{block}
\end{column} % End of column 2.1

\begin{column}{\onecolwid}\vspace{-.6in} % The second column within column 2 (column 2.2)
\begin{block}{The natural hedge mostly helps continuous corn}
		%\vspace{-0.9in}
		\begin{figure}
			\begin{center}
				\includegraphics[scale=1]{figs/fig_p_debt_all.png}\\
				\caption{\hspace{.05cm}Mean year-end rollover debt: solid = correlated prices/yields (hedge active), dashed = independent (no edge). The vertical gap between dashed and solid is the natural hedge benefit. The gap is widest for continuous corn (red) and nearly closes for the rotation plans.}
				\label{fig:fig_p_debt_all}
			\end{center}
		\end{figure} 
\end{block}
\end{column}	
\end{columns} % End of the split of column 2 - any content after this will now take up 2 columns width
\end{column} % End of the second column

\begin{column}{\sepwid}\end{column} % Empty spacer column

\begin{column}{\onecolwid} % The third column
	
\begin{block}{Rotation Plans}
\begin{table}
    \centering
    \begin{tabular}{ccc}
        Code & Pattern  & \% Prevalence\\
       \textcolor{red}{CCCCCC}  & Corn monoculture & 7\\
       \textcolor{blue}{CSCSCS} & Perfect rotation & 53\\
       \textcolor{yellow}{CCSCCS} & 2 corn/1 soy & 6 
    \end{tabular}
    \caption{CSCSCS and CCSCCS share the same two crops but differ in when soy lands\\ The lever behind both the agronomic and the price-shock results.}
    \label{tab:rot_perc}
\end{table}
\vspace{1.2cm}
\end{block}	

\begin{block}{Findings}
\begin{enumerate}
    \item Rotation lowers risk not by smoothing it but by periodically arresting debt rollover: a soy harvest retires accumulated balance, resetting the farm to near-zero delinquency before the next corn run rebuilds it.
    \item The rotation advantage compounds: fewer delinquency events plus periodic soy harvests leave more cash on the balance sheet by year 6; even though corn has the higher gross revenue per acre.
    \item Rotation and the natural hedge are partial substitutes. Continuous corn relies on the hedge because it never retires its rollover balance; rotations already reset debt to near-zero in soy years, so the hedge has little left to do.
    \item Rationing is rare at these limits, but rotation lowers peak demand; an observable, verifiable risk factor lenders can price into operating-credit underwriting.
\end{enumerate}
\end{block}

\begin{block}
\centering
 \textbf{Contact Info}\\
 Victor Funes-Leal: vf006@uark.edu\\
 Lawson Connor: lconnor@uark.edu\\
 Eunchun Park: park@uark.edu
\end{block}

\end{column} % End of the third column
\end{columns} % End of all the columns in the poster

%\begin{columns}[t]
%\begin{column}{\fourcolwid}
%\begin{textblock}{4}(12.25,-.65)
%\centering
% \textbf{Contact Info}\\
% Victor Funes: vf006@uark.edu\\
% Lawson Connor: lconnor@uark.edu\\
% Eunchun Park: park@uark.edu
%\end{textblock}
% \quad \\
%\vspace{-2cm}
%\footnotesize
%\textbf{References}\\
%\footnotesize
%Schroeder, T., Espinosa, J. and Goodwin, B. (1992). The Value of Genetic Traits in Purebred Dairy Bull Services. \emph{Review of Agricultural Economics 14} (2), 215--226.\\
%Richards, T. and Jeffrey, S. (1996) Establishing Indices of genetic Merit Using Hedonic Pricing: An Application to dairy Bulls in Alberta. \emph{Canadian Journal of Agricultural Economics 44}, 251--264. \\
%Sy, H., Faminow, M. Johnson, G. and Crow, G. (1997). Estimating the Values of Cattle Characteristics Using an Ordered Probit Model. \emph{American Journal of Agricultural Economics 79} (2), 463--476.\\
%\end{column}
%\end{columns}
\end{frame} % End of the enclosing frame
\end{document}
